\chapter{细胞内部不是装满液体的小袋子}

\begin{kaichang}
去厨房拿一枚生鸡蛋，磕开，把蛋清蛋黄完整地倒进碗里。你会看到一颗圆滚滚的蛋黄，被一层几乎看不见的薄膜包着，筷子轻轻一戳，膜破了，浓稠的黄色液体流出来。

这里藏着一个冷知识：那颗蛋黄，在生物学上算\emph{一个细胞}。包着它的蛋黄膜，就是它的细胞膜。鸡蛋黄直径约三厘米，鸵鸟蛋的蛋黄能有拳头大——它们是地球上最大的单个细胞之一。

现在请你猜一猜：刚才流出来的那摊黄色液体，就是"细胞液"吗？细胞内部，是不是就像一个装满营养汤汁的小袋子，细胞器像饺子一样漂在汤里？这是很多人对细胞的第一印象，教科书上的彩图也在强化这种印象——一个圆圈，里面漂着几块彩色的小零件。先写下你的答案。这一章结束时，我们会回到这颗蛋黄，你会发现真实细胞的内部比"一袋汤"疯狂得多。
\end{kaichang}

\section{显微镜下：细胞里从来不是空的}

先只看能看见的东西。如果你有机会用显微镜看一片洋葱表皮的临时装片，你会看到一排排像砖墙一样的小格子，每个格子里有一个颜色较深的小圆点——细胞核。把镜头换到一滴池塘水，画面立刻变成闹市：草履虫飞快地游过，硅藻像玻璃小船一样漂着，有些细胞里能清楚看到颗粒在缓缓移动。

最说明问题的是一种叫黑藻的水草。取一片嫩叶放在显微镜下，你会看到每个细胞里，绿色的叶绿体排着队、绕着细胞内壁不停地转圈，像一条环形传送带，这叫做\textbf{细胞质流动}。如果细胞内部真是一袋静止的汤，这条传送带根本转不起来——汤里没有什么东西能有方向地"推"着叶绿体跑。

再看一个你更熟悉的场景。你磕破手指，几天后伤口结痂愈合：皮肤细胞分裂、爬动、把缺口补上。细胞会\emph{爬}，会改变形状，会把物质从一头运到另一头。静止的液体袋子做不到这些事。可见光是第一件证据：细胞内部有结构，而且这些结构在不停地动。

\section{把镜头推进：浓粥，不是清汤}

继续放大，超出光学显微镜的极限，进入电子显微镜的世界。这里的第一冲击是：\textbf{细胞里挤得要命}。

给你几个数字感受一下。一个大肠杆菌（一种细菌，长约 \ud{2}{\mu m}，是你手指甲缝里可能住着几十亿个的小东西）体内大约有 $3\times10^{6}$ 个蛋白质分子，再加上几千个核糖体、一份环状 DNA 和数不清的小分子。把这些大分子全部摊开，它们要占据细胞内部大约 $30\%$ 的体积。$30\%$ 是什么概念？高峰期的地铁车厢，人均占的体积比例也就这个量级。你体内的细胞比大肠杆菌还大、还挤。

所以细胞质不是稀汤，是一锅浓到搅不动的粥。这一点至关重要，因为它决定了细胞里几乎一切事情的发生方式：一个蛋白质分子在细胞里移动，不像石头落进水里那样笔直，而是像你在春运火车站里穿行——每走一步都要和几百个"别人"碰撞、绕路。分子之间没有私人空间，酶和它的底物是被拥挤的人群"推"到一起的。你将在后面看到，细胞对付拥挤的办法不是减少人口，而是修"轨道交通"。

\section{两种基本户型：原核与真核}

地球上的细胞只有两种基本户型，差别一目了然。

\textbf{原核细胞}（细菌和古菌）：没有细胞核。DNA 盘成一团，待在细胞里一个叫"类核"的区域，但没有任何膜把它和周围隔开。除了细胞膜，内部几乎没有膜结构。它们小——通常直径一两微米——靠小分子随机扩散就能完成大部分内部运输，所以不需要复杂的内部隔间。请注意：原核细胞不等于"一锅粥"。它们的 DNA 有固定的摆放位置，蛋白质机器按分工聚集，连形状都有专门的蛋白质支架维持。只是它们不搞"带墙的办公室"。

\textbf{真核细胞}（你、蘑菇、海带、草履虫）：大得多，典型直径 $10$ 到 \ud{100}{\mu m}，内部用膜隔出了一整套"房间"，这就是\textbf{细胞器}。你的细胞、酵母的细胞、洋葱的细胞，都属于这种户型。第 56 章讲过最早的真核细胞大约出现在二十亿年前，正是这套内膜系统让细胞变大、变复杂成为可能。

先澄清一件事：膜不等于袋子。一层膜是一个\textbf{工作表面}。第 51 章讲过，磷脂双分子层对大多数带电物质不通透，这意味着膜两边可以维持不同的酸碱度、不同的离子浓度——膜本身就是一堵可以"蓄水发电"的坝。第 53 章的呼吸链、第 54 章的光合链，全都铺在膜上工作。理解了"膜是工作台"，下面这些细胞器就不再是名词清单。

\section{逐个房间看，但要看机器不看门牌}

教科书的讲法是：细胞核是"控制中心"，线粒体是"发电厂"，高尔基体是"邮局"——背完就完了。这种比喻用来记位置可以，但它完全没告诉你事情是怎样发生的。我们换个方式：每个细胞器，只问一个问题——\textbf{它的膜和蛋白质分子在干什么}。

\subsection{细胞核：有门卫的保险库}

细胞核的膜是双层的，上面开满了一种叫\textbf{核孔}的通道。核孔不是一个洞，而是一台由几十种蛋白质组成的安检机：它只放行携带特定"通行信号"（一段氨基酸序列）的蛋白质。DNA 太大，从来不离开细胞核；要执行命令时，细胞抄写一份 RNA 副本送出去（这就是第三册第 67 到 69 章要细讲的转录与翻译）。保险库加安检门的设计解决了一个真问题：真核细胞的 DNA 比细菌长上千倍，必须和细胞质里翻腾的化学反应隔离开，同时还要保持双向通信。

\subsection{内质网和高尔基体：折叠车间与分拣流水线}

内质网是一片从核膜延伸出来的、折叠成迷宫的巨大膜系统。为什么要折叠？为了在小体积里提供巨大的工作面积——和你小肠内壁的褶皱是同一个几何原理。附着核糖体的那部分叫粗面内质网：新合成的蛋白质像穿线一样直接穿进内质网腔内，在里面折叠、检查质量，折叠失败的会被标记销毁（第 49 章讲过蛋白质折叠）。不附着核糖体的光面内质网则是脂质车间，细胞膜、各种细胞器膜的材料都在这里生产。

产品怎么送到下一站？靠\textbf{囊泡}：内质网膜的局部像吹泡泡一样鼓出去、断开，变成一个装着货物的小膜泡，漂到高尔基体，膜与膜融合，货物卸货。高尔基体是一摞压扁的膜囊，像一叠盘子，货物从这头进、那头出，沿途被"贴标签"——主要是接上或修剪糖链，糖链就是蛋白质的目的地地址（第 47 章讲过糖）。贴好标签的货物再装进新囊泡，发往细胞膜、溶酶体或细胞外。

最后一个关键细节：囊泡凭什么不送错地址？答案是膜上的蛋白质互相识别。囊泡膜和目标膜上各有一组配对的蛋白质，形状互补，相遇时像拉链齿一样咬合，把两层膜拉到一起、逼它们融合。不是配对的膜，根本合不上。细胞内部的"物流"没有司机和地图，只有分子层面的锁和钥匙——这就是第 48 章蛋白质结构决定功能的一次大规模应用。

\subsection{溶酶体：用酸当保险的回收站}

溶酶体装着几十种能拆解蛋白质、脂质、糖的水解酶。一个显然的危险是：这些酶要是漏出来，细胞不就被自己消化了吗？细胞的解法很化学：溶酶体膜上有质子泵，不停消耗 ATP 把 \chem{H^+} 泵进去，把内部维持在 $pH$ 约 $5$ 的酸性环境，而这些水解酶恰恰只在酸性条件下火力全开，到了中性的细胞质里就基本熄火。酸性环境是一重保险，中性细胞质是第二重。你看，这里没有任何"智能"，只是第 34 章讲过的 $pH$ 与酶活性的关系被精确利用了。溶酶体负责回收衰老的细胞器、消化吞进来的细菌；有一类遗传病叫溶酶体贮积症，就是某种水解酶天生失灵，拆不掉的物质在细胞里越堆越多——单一一台酶失效，后果可以是致命的。

\subsection{线粒体和叶绿体：自带证件的能源部门}

线粒体和叶绿体的工作机制，第 53、54 章已经拆过：呼吸链把电子传递释放的能量变成跨膜质子梯度，质子回流推动 ATP 合酶旋转；光合链用光能完成类似的一跃。这里只强调三个在别的细胞器身上看不到的特征：它们都有\emph{双层}膜；它们体内都装着自己的 DNA 和自己的核糖体，能自己合成一小部分蛋白质；它们都靠自己一分为二地增殖，细胞没法"从零造出"一个新的线粒体——线粒体只能从已有的线粒体长出来。

一个细胞器，为什么要有自己的 DNA？这三个怪癖放在一起，指向一个惊人的故事，我们留到证据层讲。

\subsection{细胞骨架：会拆会建的脚手架和铁路}

以上所有运输和形态，都依赖一套贯穿整个细胞的蛋白质纤维网，统称\textbf{细胞骨架}。它主要有两种轨材：一种叫微管，是直径约 \ud{25}{nm} 的中空管子，由一种蛋白质零件首尾相接拼成；另一种叫微丝，更细，由肌动蛋白串成。注意"拼装"这个词——这些轨道是\emph{活的}，每一秒都在一头加装零件、另一头卸掉零件，需要时几分钟内就能在细胞里铺一条新路，用完即拆。

轨道上跑着\textbf{分子马达}。研究得最透彻的一种叫驱动蛋白：它有两只"脚"，交替着沿微管向前走，每走一步，身后就消耗一个 ATP——ATP 水解让它的形状猛地一抖，这一抖就是一步（第 50 章讲过 ATP 水解如何改变蛋白质形状）。它的步长大约 \ud{8}{nm}，每秒能走几十到上百步，背上驮着比自己大几十倍的囊泡。你此刻能思考，部分原因是你的神经细胞里，驱动蛋白正沿着长达一米的微管，把囊泡从细胞体一步步扛到神经末梢。

微丝则负责"肌肉"：两种蛋白质（肌动蛋白和肌球蛋白）互相错动，让细胞收缩、爬行、在分裂时从中间把自己勒成两半。黑藻叶绿体的环形传送带，推力正是来自贴着细胞壁铺设的微丝轨道。

\section{为什么不全靠扩散？算一笔时间账}

\begin{qiaoliang}
没有马达、没有轨道，分子靠随机碰撞（扩散）在细胞里移动，够用吗？物理学家告诉我们，扩散的平均距离随时间的\emph{平方根}增长：想走两倍远，要花四倍时间。蛋白质在细胞质里的扩散系数大约是 $D\approx\ud{10}{\mu m^2/s}$。粗略估算，分子扩散距离 $x$ 所需的时间约为 $x^2/(6D)$。

- 穿过大肠杆菌（\ud{1}{\mu m}）：约 $1/(6\times10)\approx\ud{0.02}{s}$。眨眼之间，所以细菌不需要内部运输系统。
- 穿过一个典型真核细胞（\ud{100}{\mu m}）：约 $100^2/60\approx\ud{170}{s}$，好几分钟。有点慢，所以真核细胞配了马达和轨道。
- 穿过一颗直径三厘米的鸡蛋黄（\ud{3\times10^{4}}{\mu m}）：约 $(3\times10^{4})^2/60\approx\ud{1.5\times10^{7}}{s}$，差不多\emph{半年}。

看明白了吗？扩散是"距离平方"的奴隶。细胞越大，靠随机乱走送货就越绝望。这就是为什么大细胞必须修路，以及为什么那颗巨大的蛋黄另有生存策略——回到开场时揭晓。
\end{qiaoliang}

\section{一张诚实的示意图}

下面是真核细胞的示意图和细胞器清单。请牢记：图中细胞器的大小比例、位置和形状全是人为的表示法——真实细胞里它们挤在一起、不断变形，绝不是漂在空白背景上的彩块。

\begin{center}
\begin{tikzpicture}[scale=1.0]
\draw[thick] (0,0) ellipse (5.2 and 3.4);
\draw[thick] (-1.6,0.4) ellipse (1.5 and 1.2);
\draw[thick] (-1.6,0.4) ellipse (1.2 and 0.9);
\draw (1.8,1.2) ellipse (0.9 and 0.45);
\draw (1.8,1.2) ellipse (0.6 and 0.28);
\draw (2.6,-1.2) ellipse (0.9 and 0.45);
\draw (2.6,-1.2) ellipse (0.6 and 0.28);
\draw[rounded corners=4pt] (-0.2,1.7) rectangle (1.6,2.1);
\draw[rounded corners=4pt] (-0.3,2.15) rectangle (1.5,2.55);
\draw[rounded corners=4pt] (-0.4,2.6) rectangle (1.4,3.0);
\draw (0.9,-0.6) circle (0.35);
\draw (0.2,-1.8) circle (0.28);
\draw[->,thick] (-1.6,1.7) -- (-1.6,2.9);
\draw[->,thick] (-0.2,-0.6) -- (0.55,-0.6);
\node at (-1.6,-2.1) {\small 细胞核};
\node at (4.0,1.2) {\small 线粒体};
\node at (2.2,2.7) {\small 高尔基体};
\node at (0.9,-1.2) {\small 溶酶体};
\node at (0,-3.9) {\small 囊泡沿细胞骨架轨道在细胞器之间运输（本图比例严重失真）};
\end{tikzpicture}
\end{center}

\begin{table}[htbp]
\centering
\caption{真核细胞主要细胞器速查（植物细胞特有：叶绿体、大液泡、细胞壁）}
\begin{tabular}{lll}
\toprule
细胞器 & 结构特征 & 核心机制一句话 \\
\midrule
细胞核 & 双层膜 + 核孔 & DNA 保险库，核孔蛋白选择性放行 \\
内质网 & 折叠膜迷宫 & 增大膜面积，蛋白质折叠、脂质合成 \\
高尔基体 & 膜囊堆叠 & 糖链加工贴标签，囊泡按地址分发 \\
溶酶体 & 单层膜囊 & 质子泵造酸性环境，水解酶酸性下才工作 \\
线粒体 & 双层膜，内膜成嵴 & 呼吸链与 ATP 合酶（第 53 章） \\
叶绿体 & 双层膜 + 类囊体 & 光合链（第 54 章） \\
细胞骨架 & 微管、微丝 & 拼装式轨道 + 分子马达定向运输 \\
\bottomrule
\end{tabular}
\end{table}

\section{科学家怎样知道：细胞里住着远古的房客}

\begin{zhengju}
早在 1880 年代，就有学者在显微镜下注意到叶绿体和能进行光合作用的蓝细菌长得像亲兄弟；1905 年，俄国学者梅列什科夫斯基干脆提出：叶绿体可能起源于被吞进细胞、却没被消化掉的细菌。这个想法太大胆，没DNA、没序列数据，谁也证不了，被搁置了半个多世纪。

1967 年，美国生物学家林恩·马古利斯把"线粒体和叶绿体起源于被吞并的细菌"这个\textbf{内共生理论}整理成论文投出去，结果被十几家期刊退稿，理由是证据不足、异想天开。她坚持把论文发了出来，然后和整个生物学界一起，等了十几年——等的是分子生物学的新工具。新证据一条条落下，而且每一条都直指细菌的"身份证"：

\begin{itemize}[itemsep=2pt]
\item \textbf{DNA 证据}：线粒体和叶绿体里有自己的 DNA，而且和细菌一样是环状的。更致命的一击来自序列比对（第 66 章以后你会看到原理）：把线粒体 DNA 的序列和各类生物逐一比较，它和一类叫 $\alpha$-变形菌的细菌最亲近；叶绿体的序列则明确指向蓝细菌。如果它们只是细胞自己长出来的膜结构，完全没道理自带一套"细菌式样"的遗传物质。
\item \textbf{膜证据}：它们都有双层膜，而且内膜的脂质成分像细菌细胞膜，外膜则像宿主细胞吞食物时的吞噬泡膜——一次"吞了没消化"的事件，正好留下这样两层来历不同的膜。它们还像细菌一样靠一分为二地分裂增殖，细胞无法凭空制造新的。
\item \textbf{核糖体证据}：线粒体和叶绿体内部的核糖体是"细菌型号"，而且会被某些专攻细菌核糖体的抗生素抑制；同一个细胞里，细胞质中的核糖体却对这些抗生素不敏感。一套房子里两套不同型号的蛋白质制造机，用"同宗不同源"才解释得通。
\end{itemize}

替代解释也被认真检验过：会不会这些相似纯属巧合，或者是细胞核基因"搬"进去的结果？序列证据否决了它们——相似不是表面像，而是逐字母的像。今天，科学家甚至能观察到"进行中"的内共生：某些单细胞生物体内住着刚被收编不久的蓝细菌，正处在房客变器官的过渡阶段。这个故事值得记住的不是马古利斯多有远见，而是：一个被退了十几次稿的理论，只要它做出了可以被检验的预言，就等来了翻盘。顺便说，你的线粒体 DNA 全部来自母亲的卵细胞——父亲精子里的线粒体在受精后会被清除，这给第三册的遗传故事留了个伏笔。
\end{zhengju}

\begin{shiyan}
\textbf{看一看细胞在你碗里干活}（安全，可在家做）

材料：一小袋干酵母（超市有售）、白糖、温水（手摸不烫，约 $35\,^{\circ}\mathrm{C}$）、两个透明瓶子、两个气球。

步骤：两瓶各装半瓶温水、各加一勺糖；一瓶加半勺干酵母，一瓶不加，摇匀；瓶口各套一个气球，放在温暖处，过半小时到一小时观察。

观察点：加了酵母的瓶子里，气球会慢慢鼓起来，液面冒细密的气泡；没加酵母的那瓶毫无动静。气球里的气体主要是 \chem{CO_2}——酵母细胞正用第 52 章的糖酵解拆解糖，边拆边排气。你看不到细胞内部，但整个瓶子被酵母"改了化学"，这就是细胞不是空袋子的宏观证据。

安全提示：全程只用食品级材料，无毒无害；不要用热水，会烫死酵母也让实验失去对照；实验后液体倒进水槽冲走即可。
\end{shiyan}

\section{本章模型在哪里会失灵}

第一，教科书式的"细胞器分布图"全是定格照片。真实的线粒体不是一颗颗独立的小香肠，而是一张不断融合、分裂、重构的动态网络，几秒钟就变一次拓扑；内质网是连续的一整片；高尔基体在细胞分裂时会先解体再重建。把细胞器当成固定的"器官零件"，是这个模型最常用的错法。

第二，"原核细胞结构简单"不等于"原核细胞是一锅匀汤"。细菌的 DNA 有确定的折叠和定位，维持杆状外形的骨架蛋白（和真核细胞骨架同宗）按几何规则排列。只是它们的组织方式不依赖膜隔间。

第三，"真核 = 有核、原核 = 无核"在已知生物里很好用，但边界正在松动：科学家在深海沉积物里发现的一些古菌，拥有不少原本认为真核生物独有的基因，它们可能是真核细胞那位远古"房东"的亲戚。内共生和细胞起源的细节——宿主到底是谁、吞噬先后如何——至今仍在争论和修订中，这正是科学还活着的样子。

\begin{wuqu}
"植物细胞有叶绿体，所以不需要线粒体。"——这是把光合作用和呼吸作用当成你死我活的替代品了。事实是：叶绿体只在有光时、只在绿色部分工作，把光能转成糖里的化学能；而植物每一个活细胞——包括永不见光的根细胞、包括深夜里的叶片细胞——都要靠线粒体把糖的能量转成可直接使用的 ATP（第 53 章）。一颗埋在土里的种子，发芽时没有一片绿叶，全靠线粒体供能。反例就在你家窗台：把一盆绿植关进暗室几天，它不会立刻死去，因为线粒体一直在值班；但没有线粒体的植物细胞，一刻也活不下去。
\end{wuqu}

\section{回到那颗蛋黄}

现在可以对答案了。磕开鸡蛋流出来的蛋黄液，与其说是"细胞液"，不如说是一间巨大的\emph{补给仓库}：蛋黄的绝大部分体积被卵黄蛋白、脂质等储备物资占据，是未来胚胎二十一天的口粮；真正拥挤、忙碌、布满骨架和马达的"活"细胞质，只是蛋黄表面一个直径几毫米的小白点——胚盘。所有细胞分裂、分化、搭建小鸡身体的工程，都在那个小点及其边缘进行。

为什么这样安排？桥梁层那笔时间账就是答案：一颗直径三厘米的细胞如果处处靠扩散送货，横穿要半年，谁都等不起。卵细胞的演化对策是：把"仓库"做大，把"工厂"做小。所以"细胞是装满液体的小袋子"这个印象，连世界上最大的细胞自己都不买账——袋子里装的不是汤，是弹药；而干活的那个角落，拥挤、喧闹、有铁路、有马达，还有住着几十亿年的远古房客。

\section{继续深挖：铁路成了药，消息怎么进门}

理解了细胞骨架，就能理解一类抗癌药：紫杉醇最初从红豆杉树皮中提取，它恰好卡住微管"边装边拆"的动态平衡，把轨道冻住。细胞分裂时染色体要靠微管牵引，轨道一冻，快速分裂的癌细胞就被卡在半路口死去（正常细胞也受波及，这就是化疗副作用的来源，第 59 章还会从细胞合作的角度再看癌症）。

还有一个更基本的问题被我们整个跳过了：内质网接单、高尔基体发货、溶酶体调酸、马达发车——这些调度由谁下令？细胞又怎么知道外面来了养分、危险或邻居的信号？膜上有一整套"天线"和"传令链"，能把膜外的一个分子信号，放大成细胞内部的千军万马。下一章，第 58 章，我们就去看细胞怎样听见外界的消息。

\begin{tiaozhan}
扩散为什么那么慢？把分子的随机乱走想成抛硬币：每走一步，方向随机。数学上可以证明，走 $N$ 步之后，离起点的均方距离正比于 $N$——也就是说，距离正比于步数的平方根。换成连续语言：$\langle x^2\rangle = 6Dt$（三维情形），时间 $t$ 与距离的\emph{平方}成正比。这正是桥梁层那笔账的来源，也是"平方根"这个数学形状反复出现在热学、布朗运动、甚至股票波动里的原因。另一个厉害的估算来自几何：细胞体积（和它的食量）随直径的立方增长，膜面积（和它的进食能力）只随直径的平方增长——所以细胞直径翻倍，每个细胞器要"喂饱"的体积翻八倍。这个立方对平方的赛跑，是细胞为什么长不大、卵细胞为什么必须特殊设计的几何根源。跳过本框不影响主线。
\end{tiaozhan}

\section*{练一练}

\begin{sikaoti}
\item 画一幅"蛋白质出厂物流图"：从核糖体开始，用箭头画出一件膜蛋白依次经过内质网、囊泡、高尔基体、囊泡、细胞膜的路径，并在每个站点旁边写一句"这里的分子在干什么"。
\item 参照桥梁层的算法，估算一个物质靠扩散穿过直径 \ud{10}{\mu m} 的细胞核大约要多久（取 $D\approx\ud{10}{\mu m^2/s}$），再算出穿过一米长的神经细胞轴突要多久，由此说明为什么轴突里必须有分子马达。
\item 内共生理论说线粒体起源于被吞并的细菌。请写出本章给出的三类证据（DNA、膜、核糖体），并为每一类补一句：如果线粒体是细胞自己造出来的膜结构，预期观察会是什么、实际观察又是什么？
\item 溶酶体用"内部酸性、酶只在酸性下高效工作"做保险。请从第 34 章 $pH$ 影响酶活性的角度，画出你设想的溶酶体水解酶活性随 $pH$ 变化的曲线草图，并标出溶酶体内和细胞质中的两个工作点。
\item 一个原核细胞没有线粒体，但它同样要合成 ATP。根据第 53 章"呼吸链铺在膜上"的原理，推测细菌的呼吸链可能铺在哪里？这个推测和内共生理论关于线粒体内膜来历的说法，是否互相印证？
\end{sikaoti}
